\section{Distributed and minimal localization}
\label{sec:distributed}

A phenomenon need not be available at any single receiver. Once realization is phrased as a factorization problem, the same criterion applies to collections of receivers without introducing an attribution score or a new notion of sufficiency.

For $K\subseteq I_{q+1}$, define the joint exposure
\[
Q_{q,K}^{\theta,C}
:=
\langle Q_{q,j}^{\theta,C}\rangle_{j\in K}
\]
corestricted to its realized image. For $K=\varnothing$, take the unique map to a singleton. Two supported states have the same joint exposure exactly when every receiver in $K$ sees them as the same state, so
\begin{equation}
\Ker Q_{q,K}^{\theta,C}
=
\bigcap_{j\in K}D_{q,j}^{\theta,C},
\label{eq:joint-kernel}
\end{equation}
where the empty intersection is the universal relation on $S_q^{\theta,C}$.

Combining \eqref{eq:joint-kernel} with \cref{thm:realization-criterion}, the receiver family $K$ realizes $\Phi$ exactly when
\begin{equation}
\bigcap_{j\in K}D_{q,j}^{\theta,C}
\subseteq
D_{\Phi,q}^{\theta,C}.
\label{eq:family-realization}
\end{equation}
This expression makes the role of a distributed family explicit. Each receiver removes some distinctions. The family succeeds when their joint state separates every pair that the phenomenon needs separated.

The joint next-layer exposure factors the complete update. Let
\[
Q_q^{\theta,C}:=Q_{q,I_{q+1}}^{\theta,C}.
\]
On its realized joint image, the receiver continuations assemble into a map $G_q^{\theta,C}$ with
\begin{equation}
F_q^{\theta,C}=G_q^{\theta,C}Q_q^{\theta,C}.
\label{eq:joint-factorization}
\end{equation}
If $q<r$ and $\Phi_q^{\theta,C}$ is the pullback of a profile at stage $r$, then
\[
\Phi_q^{\theta,C}
=
\Phi_r^{\theta,C}
F_{q+1\rightsquigarrow r}^{\theta,C}
G_q^{\theta,C}Q_q^{\theta,C}.
\]
Hence
\begin{equation}
\Ker Q_q^{\theta,C}
\subseteq
\Ker\Phi_q^{\theta,C}.
\label{eq:full-joint-realization}
\end{equation}
The declared next-layer receiver organization jointly retains every distinction required for a later phenomenon, although no single receiver need do so.

\begin{definition}[Distributed receiver realization]
A phenomenon is \emph{distributed at cut $q$ relative to the declared receiver decomposition} if no individual receiver realizes it but some family $K$ with $|K|>1$ satisfies \eqref{eq:family-realization}.
\end{definition}

For example, on $S=\{0,1\}^2$ with
\[
\Phi(a,b)=a\oplus b,
\qquad
Q_1(a,b)=a,
\qquad
Q_2(a,b)=b,
\]
neither receiver realizes parity individually, while the joint exposure $(a,b)$ does.

\subsection{Realizing families and localization order}

Define the collection of realizing receiver families by
\begin{equation}
\Rcal_{\Phi,q}^{\theta,C}
:=
\left\{
K\subseteq I_{q+1}:
\bigcap_{j\in K}D_{q,j}^{\theta,C}
\subseteq
D_{\Phi,q}^{\theta,C}
\right\}.
\label{eq:realizing-families}
\end{equation}
It is upward closed: once $K$ realizes the phenomenon, adding receivers can only refine the joint fibers.

The minimum receiver count is
\begin{equation}
d_{\Phi,q}^{\theta,C}
:=
\begin{cases}
\min\{|K|:K\in\Rcal_{\Phi,q}^{\theta,C}\},
&\Rcal_{\Phi,q}^{\theta,C}\neq\varnothing,\\[0.3em]
+\infty,&\text{otherwise.}
\end{cases}
\label{eq:realization-order}
\end{equation}
For a nonconstant downstream profile with $q<r$, \eqref{eq:full-joint-realization} gives
\[
1\le d_{\Phi,q}^{\theta,C}\le |I_{q+1}|.
\]
A constant profile has order $0$.

The same condition has an equivalent pair-cover form. Define the phenomenon-distinguished pairs
\[
P_{\Phi,q}^{\theta,C}
:=
\left\{
(x,x')\in (S_q^{\theta,C})^2:
\Phi_q^{\theta,C}(x)\neq\Phi_q^{\theta,C}(x')
\right\},
\]
and for each receiver $j$ let
\[
E_{q,j}^{\theta,C}
:=
\left\{
(x,x')\in P_{\Phi,q}^{\theta,C}:
Q_{q,j}^{\theta,C}(x)\neq Q_{q,j}^{\theta,C}(x')
\right\}.
\]

\begin{proposition}[Pair-cover characterization]
\label{prop:pair-cover}
A receiver family $K\subseteq I_{q+1}$ realizes $\Phi$ if and only if
\begin{equation}
P_{\Phi,q}^{\theta,C}
\subseteq
\bigcup_{j\in K}E_{q,j}^{\theta,C}.
\label{eq:pair-cover}
\end{equation}
Consequently,
\begin{equation}
d_{\Phi,q}^{\theta,C}
=
\min\left\{
|K|:
K\subseteq I_{q+1},\ 
P_{\Phi,q}^{\theta,C}
\subseteq
\bigcup_{j\in K}E_{q,j}^{\theta,C}
\right\},
\label{eq:pair-cover-order}
\end{equation}
where the minimum is $+\infty$ if no such $K$ exists. The empty family realizes exactly when $\Phi_q^{\theta,C}$ is constant.
\end{proposition}

\begin{proof}
By \eqref{eq:family-realization}, $K$ fails to realize $\Phi$ exactly when there is a pair $(x,x')$ distinguished by $\Phi_q^{\theta,C}$ that lies in every receiver kernel $D_{q,j}^{\theta,C}$ for $j\in K$. Equivalently, $K$ realizes $\Phi$ exactly when every pair in $P_{\Phi,q}^{\theta,C}$ is separated by at least one receiver in $K$, which is \eqref{eq:pair-cover}. The formula for $d_{\Phi,q}^{\theta,C}$ follows from \eqref{eq:realization-order}.
\end{proof}

Thus the localization order is the minimum number of declared receivers needed to cover all state-pair distinctions demanded by the phenomenon.

The count $d_{\Phi,q}^{\theta,C}$ is useful but incomplete. A circuit claim normally identifies components, not only their number. We therefore also care about the inclusion-minimal elements of $\Rcal_{\Phi,q}^{\theta,C}$. The next section shows that both the minimum count and the identity of these minimal families depend on support.
